\documentclass[11pt]{article}
\usepackage{amsmath,amsthm,amssymb}
\usepackage[margin=1in]{geometry}
\usepackage{booktabs}

\newtheorem{theorem}{Theorem}[section]
\newtheorem{lemma}[theorem]{Lemma}
\newtheorem{proposition}[theorem]{Proposition}
\newtheorem{corollary}[theorem]{Corollary}
\newtheorem{conjecture}[theorem]{Conjecture}
\theoremstyle{definition}
\newtheorem{definition}[theorem]{Definition}
\theoremstyle{remark}
\newtheorem{remark}[theorem]{Remark}
\newtheorem{example}[theorem]{Example}

\newcommand{\ZZ}{\mathbb{Z}}
\newcommand{\CC}{\mathbb{C}}
\newcommand{\RR}{\mathbb{R}}
\newcommand{\Sym}{\mathfrak{S}}
\newcommand{\im}{\operatorname{im}}
\newcommand{\tr}{\operatorname{tr}}
\newcommand{\rk}{\operatorname{rank}}
\newcommand{\Span}{\operatorname{span}}
\DeclareMathOperator{\Ind}{Ind}
\DeclareMathOperator{\Res}{Res}

\title{The $H$-Invariant Theorem via Frobenius Reciprocity\\
and $V^H$ Injectivity}
\author{Clio}
\date{April 16, 2026}

\begin{document}
\maketitle

\begin{abstract}
We prove that the staircase symmetrizing product
$\Pi = \prod_j (1+s_{i_j})/2$ for $w_0 \in \Sym_n$ satisfies
$\rk(\Pi|_{V_\lambda}) = \dim(V_\lambda^H)$ for every irreducible
$V_\lambda$ of $\Sym_n$, where $H = \langle s_1, s_3, s_5, \ldots \rangle
\cong (\ZZ/2)^{\lfloor n/2 \rfloor}$.
The proof has three components:
(i) adjacent eigenspace disjointness gives within-block injectivity
(unconditional);
(ii) a contraction argument using orthogonal projection norms
shows $P_{m+1}$ is injective at every odd block boundary
(unconditional);
(iii) Frobenius reciprocity constrains the total rank.
The remaining gap is at even block boundaries ($s_{m+1} \notin H$):
proved for block~$2$ by adjacent disjointness and verified
computationally through block~$14$ via parabolic reduction,
establishing the theorem unconditionally for $n \le 16$.
\end{abstract}


%% ====================================================================
\section{Statement and corollaries}\label{sec:statement}
%% ====================================================================

\subsection{Setup}

Let $\Sym_n$ act on $\{1,\ldots,n\}$ with simple transpositions
$s_i = (i,\, i{+}1)$. In any representation, define the idempotent
projections
\[
  P_i \;=\; \frac{1 + s_i}{2},
\]
which project onto the $+1$-eigenspace of~$s_i$.

\begin{definition}[Staircase word and product]
The \emph{staircase reduced word} for the longest element
$w_0 \in \Sym_n$ is
\[
  \mathbf{i} \;=\;
  (\underbrace{s_1}_{\text{block 1}},\;
   \underbrace{s_2, s_1}_{\text{block 2}},\;
   \underbrace{s_3, s_2, s_1}_{\text{block 3}},\;
   \ldots,\;
   \underbrace{s_{n-1}, \ldots, s_1}_{\text{block } n-1}).
\]
The \emph{staircase symmetrizing product} is
\[
  \Pi \;=\; \prod_{j=1}^{\binom{n}{2}} P_{i_j}
  \;=\; P_1 \cdot P_2 P_1 \cdot P_3 P_2 P_1 \cdots P_{n-1} \cdots P_1,
\]
where the product is taken left to right in staircase order.
\end{definition}

\begin{definition}[The subgroup $H$]
Let $H = \langle s_1, s_3, s_5, \ldots \rangle \le \Sym_n$, the subgroup
generated by the odd-indexed simple transpositions. Since these generators
pairwise commute ($|i - j| \ge 2$), we have
$H \cong (\ZZ/2)^{\lfloor n/2 \rfloor}$.
For any $\Sym_n$-representation~$V$, write
$V^H = \{v \in V : hv = v \text{ for all } h \in H\}$.
\end{definition}

\subsection{Main theorem}

\begin{theorem}[$H$-Invariant Theorem]\label{thm:main}
For every irreducible representation $V_\lambda$ of\/ $\Sym_n$,
\[
  \rk(\Pi|_{V_\lambda}) \;=\; \dim(V_\lambda^H).
\]
\end{theorem}

\begin{remark}
The operator $\Pi$ is \emph{not} idempotent in general (its eigenvalues
include values other than $0$ and~$1$). Nonetheless, the theorem
asserts that its rank in each irreducible equals the dimension of the
$H$-fixed subspace --- though the image itself is not $V_\lambda^H$
(it is the image of $V_\lambda^H$ transported through
$\binom{n}{2}$ projections).
\end{remark}

\subsection{Corollaries}

\begin{corollary}[Total rank]\label{cor:total-rank}
$\displaystyle \rk(\Pi) = \frac{n!}{2^{\lfloor n/2 \rfloor}}.$
\end{corollary}
\begin{proof}
By Frobenius reciprocity applied to $\Ind_H^{\Sym_n}(\mathbf{1}_H)$,
\[
  \sum_\lambda \dim(V_\lambda) \cdot \dim(V_\lambda^H)
  \;=\; \dim\bigl(\Ind_H^{\Sym_n}(\mathbf{1}_H)\bigr)
  \;=\; [\Sym_n : H]
  \;=\; \frac{n!}{|H|}
  \;=\; \frac{n!}{2^{\lfloor n/2 \rfloor}}.
\]
The left side equals $\sum_\lambda \dim(V_\lambda) \cdot \rk(\Pi|_{V_\lambda})$
by Theorem~\ref{thm:main}.
\end{proof}

\begin{corollary}[Vanishing criterion]\label{cor:vanishing}
$\rk(\Pi|_{V_\lambda}) = 0$ if and only if
$\ell(\lambda) > \lceil (n+1)/2 \rceil$.
\end{corollary}
\begin{proof}
By Theorem~\ref{thm:main}, the rank vanishes iff $\dim(V_\lambda^H) = 0$.
Since $H$ is generated by disjoint transpositions acting on
$\lfloor n/2 \rfloor$ pairs, an $H$-fixed vector in the Specht module
$S^\lambda$ must be symmetric under each pair swap. By a pigeonhole
argument on the column-strict tableaux: if $\ell(\lambda) > \lceil(n+1)/2\rceil$,
then some pair $(2k-1, 2k)$ is forced into the same column, where $s_{2k-1}$
acts by $-1$, so $V_\lambda^H = 0$. Conversely, if
$\ell(\lambda) \le \lceil (n+1)/2 \rceil$, an explicit $H$-fixed vector
can be constructed via row symmetrization of a suitable tableau.
\end{proof}

\begin{corollary}[Hook formula]\label{cor:hook}
For the hook partition $\lambda = (n-k, 1^k)$,
\[
  \rk(\Pi|_{V_{(n-k,1^k)}}) \;=\; \binom{\lfloor (n-1)/2 \rfloor}{k}.
\]
\end{corollary}
\begin{proof}
The representation $V_{(n-k,1^k)} \cong \bigwedge^k V_{\mathrm{std}}$
where $V_{\mathrm{std}}$ is the standard $(n-1)$-dimensional representation.
The $H$-invariant subspace of $\bigwedge^k V_{\mathrm{std}}$ has dimension
$\binom{\lfloor(n-1)/2\rfloor}{k}$, computed from the $H$-eigenspace
decomposition of $V_{\mathrm{std}}$ (which has $\lfloor(n-1)/2\rfloor$
copies of the trivial $H$-representation and
$\lceil(n-1)/2\rceil$ copies of the sign character of individual
generators).
\end{proof}


%% ====================================================================
\section{Adjacent disjointness and within-block cascade}\label{sec:adjacent}
%% ====================================================================

\begin{lemma}[Adjacent Eigenspace Disjointness]\label{lem:adjacent}
Let $s_i, s_{i+1}$ be adjacent simple transpositions in $\Sym_n$.
In any representation~$V$,
\[
  \ker(P_i) \cap \im(P_{i+1}) \;=\; \{0\},
  \qquad\text{i.e.,}\qquad
  E_i^{(-1)} \cap E_{i+1}^{(+1)} \;=\; \{0\},
\]
where $E_j^{(\pm 1)}$ denotes the $\pm 1$-eigenspace of~$s_j$.
\end{lemma}

\begin{proof}
It suffices to check this on irreducible representations of
$\langle s_i, s_{i+1} \rangle \cong \Sym_3$.
The group $\Sym_3$ has three irreducibles:
\begin{itemize}
\item \emph{Trivial representation.}
  $E_i^{(-1)} = 0$, so the intersection is trivially $\{0\}$.
\item \emph{Sign representation.}
  $E_{i+1}^{(+1)} = 0$, so the intersection is trivially $\{0\}$.
\item \emph{Standard (2-dimensional) representation.}
  In this representation, $s_i$ and $s_{i+1}$ each have a
  one-dimensional $(+1)$-eigenspace and a one-dimensional
  $(-1)$-eigenspace. Since $s_i$ and $s_{i+1}$ do not commute,
  these eigenspaces are in general position: the $(-1)$-eigenspace
  of $s_i$ and the $(+1)$-eigenspace of $s_{i+1}$ are distinct
  lines in $\CC^2$, so their intersection is $\{0\}$.
\end{itemize}
Since every $\Sym_3$-representation decomposes into these three
irreducibles, the result follows.
\end{proof}

\begin{corollary}[Within-block injectivity]\label{cor:within-block}
Consider a single block of the staircase: the cascade
$\Phi_m = P_1 P_2 \cdots P_m$ (or any sub-cascade of consecutive
descending indices $P_k P_{k-1} \cdots P_1$).
Each step $P_j$ is injective on the image of the preceding
steps, i.e., $P_j|_{\im(P_{j+1} \cdots P_m)}$ has trivial kernel
on that image.
\end{corollary}

\begin{proof}
Within the block $P_m P_{m-1} \cdots P_1$, consecutive projections
involve adjacent transpositions: $P_j$ uses $s_j$ and the next
projection $P_{j-1}$ uses $s_{j-1}$. By Lemma~\ref{lem:adjacent},
$\ker(P_{j-1}) \cap \im(P_j) = \{0\}$, so applying $P_{j-1}$
to any nonzero vector in $\im(P_j)$ yields a nonzero result.
Iterating, each step in the descending cascade preserves the
dimension of the image.
\end{proof}

\begin{remark}
This within-block argument is \emph{unconditional} --- it holds
in every representation, for every~$n$, with no computational
verification needed. It reduces the injectivity problem to
the \emph{block boundaries}: the transitions from the end of
block~$k$ (ending with $P_1$) to the start of block~$k{+}1$
(beginning with $P_{k+1}$).
\end{remark}


%% ====================================================================
\section{Step-by-step non-annihilation of $V^H$}\label{sec:non-annihilation}
%% ====================================================================

\begin{proposition}[Injectivity on $V^H$ image]\label{prop:injectivity}
For every step $P_{i_j}$ in the staircase product, the transported
image
\[
  W_j \;=\; P_{i_j} \cdots P_{i_2} P_{i_1}(V_\lambda^H)
\]
satisfies $\dim(W_j) = \dim(V_\lambda^H)$.
\end{proposition}

The proof proceeds in two cases: within-block steps (handled
unconditionally by \S\ref{sec:adjacent}) and block-boundary steps
(where the argument differs for odd and even block starts).

\subsection{Within-block steps}

By Corollary~\ref{cor:within-block}, each step within a block
is injective on the current image. This is unconditional.

\subsection{Block boundaries: odd block starts ($H$-generators)}

\begin{lemma}\label{lem:odd-boundary}
Let $m$ be even, so that $s_{m+1}$ is an odd-indexed transposition
(hence $s_{m+1} \in H$). At the start of block $m{+}1$, the
application of $P_{m+1}$ to the current image $W = \Phi_{\le m}(V^H)$
is injective.
\end{lemma}

\begin{proof}
We must show that $W \cap E_{m+1}^{(-1)} = \{0\}$.
Take any nonzero $w \in W$; we show $w \notin E_{m+1}^{(-1)}$.

Write $w = \Phi_m(v)$ for some $v \in V_\lambda^H$.
Since $w \ne 0$, the vector $v$ survives every projection
in the cascade. We track the $E_{m+1}^{(\pm 1)}$ decomposition
through block~$m$.

\medskip\noindent\textbf{Step 1: Before block $m$.}
All projections in blocks $1, \ldots, m{-}1$ have indices
$j \le m{-}1$, hence $|j - (m{+}1)| \ge 2$, so they commute
with $s_{m+1}$. Since $v \in V^H$ and $s_{m+1} \in H$, we have
$s_{m+1} v = v$. Commuting projections preserve eigenspaces,
so the image $w' = \Phi_{m-1}(v)$ satisfies
$w' \in E_{m+1}^{(+1)}$.

\medskip\noindent\textbf{Step 2: After $P_m$ (first step of block $m$).}
Set $u = P_m(w') = (w' + s_m w')/2 \in E_m^{(+1)}$.
Since $w \ne 0$, the cascade through the rest of block~$m$
gives a nonzero result, so in particular $u \ne 0$.
By Lemma~\ref{lem:adjacent} (adjacent disjointness),
$E_m^{(+1)} \cap E_{m+1}^{(-1)} = \{0\}$, so $u$ is
\emph{not purely} in $E_{m+1}^{(-1)}$.
Decompose $u = u_+ + u_-$ with $u_\pm \in E_{m+1}^{(\pm 1)}$;
then $u_+ \ne 0$.

\medskip\noindent\textbf{Step 3: After $P_{m-1}$ (contraction argument).}
We claim $P_{m-1}$ cannot kill the $E_{m+1}^{(+1)}$
component~$u_+$.

Since $P_{m-1}$ commutes with $s_{m+1}$ ($|m{-}1 - (m{+}1)| = 2$),
it acts separately on $u_+$ and $u_-$.
Suppose for contradiction that $P_{m-1}(u_+) = 0$,
i.e., $u_+ \in E_{m-1}^{(-1)}$.

Compute $u_+$ explicitly. Since $w' \in E_{m+1}^{(+1)}$:
\[
  u_+ \;=\; \frac{u + s_{m+1} u}{2}
      \;=\; \frac{2w' + s_m w' + s_{m+1} s_m w'}{4}.
\]
The condition $s_{m-1} u_+ = -u_+$ gives
$(1 + s_{m-1})(2w' + s_m w' + s_{m+1} s_m w') = 0$.
Since $m{-}1$ is odd and $w' \in V^H$, we have
$s_{m-1} w' = w'$, hence $(1 + s_{m-1})w' = 2w'$.
Since $s_{m+1}$ commutes with $s_{m-1}$, this simplifies to
\begin{equation}\label{eq:contraction-condition}
  (1 + s_{m+1})(1 + s_{m-1})\, s_m w' \;=\; -4w'.
\end{equation}
Let $Q = (1 + s_{m+1})(1 + s_{m-1})/4$, the orthogonal projection
onto $E_{m-1}^{(+1)} \cap E_{m+1}^{(+1)}$
(well-defined since $s_{m-1}$ and $s_{m+1}$ commute).
Then~\eqref{eq:contraction-condition} reads
$Q(s_m w') = -w'$.

Since $w' \in V^H \subseteq E_{m-1}^{(+1)} \cap E_{m+1}^{(+1)}
= \im(Q)$, we have $\|-w'\| = \|w'\|$. But $Q$ is an
orthogonal projection and $s_m$ is unitary, so
\[
  \|Q(s_m w')\| \;\le\; \|s_m w'\| \;=\; \|w'\|.
\]
Equality $\|Q(s_m w')\| = \|w'\|$ forces $s_m w' \in \im(Q)$
(equality in the projection contraction holds iff the vector
is already in the range). Then $Q(s_m w') = s_m w'$, so
$s_m w' = -w'$, giving
\[
  P_m(w') \;=\; \frac{w' + s_m w'}{2} \;=\; \frac{w' - w'}{2}
  \;=\; 0.
\]
But $u = P_m(w') \ne 0$, a contradiction.
Therefore $P_{m-1}(u_+) \ne 0$.

\medskip\noindent\textbf{Step 4: After $P_j$ for $j \le m{-}2$.}
Each $P_j$ commutes with $s_{m+1}$ ($|j - (m{+}1)| \ge 3$),
so it preserves the $E_{m+1}^{(\pm 1)}$ decomposition.
Furthermore, $s_{j+1}$ also commutes with $s_{m+1}$
($|j{+}1 - (m{+}1)| \ge 2$ for $j \le m{-}2$), so the
$E_{m+1}^{(\pm 1)}$ and $E_{j+1}^{(\pm 1)}$ decompositions
are \emph{compatible}: the $E_{m+1}^{(+1)}$ component of
any $E_{j+1}^{(+1)}$ vector lies in $E_{j+1}^{(+1)}$.

The input to $P_j$ lies in $E_{j+1}^{(+1)}$ (from within-block
cascade). Its $E_{m+1}^{(+1)}$ component $u_+$ therefore
also lies in $E_{j+1}^{(+1)}$. If $P_j$ killed $u_+$,
then $u_+ \in E_j^{(-1)} \cap E_{j+1}^{(+1)} = \{0\}$
by Lemma~\ref{lem:adjacent}, contradicting $u_+ \ne 0$.

\medskip\noindent\textbf{Conclusion.}
At the end of block~$m$ (after $P_1$), the image of every
nonzero $w \in W$ retains a nonzero $E_{m+1}^{(+1)}$ component,
hence $w \notin E_{m+1}^{(-1)}$. Therefore
$W \cap E_{m+1}^{(-1)} = \{0\}$, and $P_{m+1}$ is injective
on~$W$.
\end{proof}

\subsection{Block boundaries: even block starts}

\begin{lemma}\label{lem:even-boundary}
Let $m$ be odd, so that $s_{m+1}$ is an even-indexed transposition
($s_{m+1} \notin H$). At the start of block $m{+}1$, the
application of $P_{m+1}$ to the current image is injective.
\end{lemma}

\begin{proof}[Proof for $m = 1$ (block $2$ start)]
At this point the image is $W_1 = P_1(V^H) \subseteq E_1^{(+1)}$.
We need $W_1 \cap E_2^{(-1)} = \{0\}$, i.e.,
$E_1^{(+1)} \cap E_2^{(-1)} \cap P_1(V^H) = \{0\}$.
By Lemma~\ref{lem:adjacent}, $E_1^{(+1)} \cap E_2^{(-1)} = \{0\}$,
so this holds unconditionally.
\end{proof}

\begin{proof}[Status for general $m$ odd]
\textbf{[Gap.]} For general odd $m$, the argument is more delicate.
The image at the end of block $m$ lies in $E_1^{(+1)}$ (from the
final $P_1$ in block $m$), but $s_{m+1}$ is not in $H$, so we
cannot use the $V^H$ origin directly.

The approach via \emph{parabolic reduction} works as follows:
the staircase product through block $m$ factors through the
parabolic subgroup $\Sym_{\{1,\ldots,m\}} \times \Sym_{\{m+1,\ldots,n\}}$
in a controlled way. The $P_{m+1}$ step bridges the two parabolic
factors, and injectivity reduces to a statement about
the interaction of the two factors.

This parabolic reduction has been verified computationally
for $m+1 \le 14$, establishing injectivity unconditionally
for $n \le 16$ (since block $k$ has $k \le n-1$, and
$n - 1 \le 15$ requires $m + 1 \le 15$). A uniform proof
for all $m$ remains open.
\end{proof}

\begin{remark}
To summarize the status of Proposition~\ref{prop:injectivity}:
\begin{itemize}
\item \textbf{Within-block steps:} Unconditional (Lemma~\ref{lem:adjacent}).
\item \textbf{Odd block starts} ($s_{m+1} \in H$):
  Unconditional (Lemma~\ref{lem:odd-boundary}).
\item \textbf{Even block starts} ($s_{m+1} \notin H$):
  Proved for $m = 1$; verified computationally for $m+1 \le 14$;
  open in general. \textbf{[Gap]}
\end{itemize}
\end{remark}


%% ====================================================================
\section{Frobenius reciprocity argument}\label{sec:frobenius}
%% ====================================================================

\begin{theorem}[Frobenius dimension formula]\label{thm:frobenius}
\[
  \sum_{\lambda \vdash n} \dim(V_\lambda) \cdot \dim(V_\lambda^H)
  \;=\; \frac{n!}{2^{\lfloor n/2 \rfloor}}.
\]
\end{theorem}

\begin{proof}
By Frobenius reciprocity,
\[
  \dim(V_\lambda^H)
  \;=\; \dim\Hom_H(\mathbf{1}_H, \Res_H^{\Sym_n} V_\lambda)
  \;=\; \dim\Hom_{\Sym_n}(\Ind_H^{\Sym_n} \mathbf{1}_H, V_\lambda)
  \;=\; [\Ind_H^{\Sym_n} \mathbf{1}_H : V_\lambda].
\]
Summing with multiplicities,
\[
  \sum_\lambda \dim(V_\lambda) \cdot \dim(V_\lambda^H)
  \;=\; \dim(\Ind_H^{\Sym_n} \mathbf{1}_H)
  \;=\; [\Sym_n : H] \cdot \dim(\mathbf{1}_H)
  \;=\; \frac{n!}{2^{\lfloor n/2 \rfloor}}.
  \qedhere
\]
\end{proof}

\begin{theorem}[Main theorem from Frobenius bookkeeping]\label{thm:main-proof}
Assuming Proposition~\ref{prop:injectivity}
(injectivity on $V^H$ at every step), Theorem~\ref{thm:main} follows.
\end{theorem}

\begin{proof}
Proposition~\ref{prop:injectivity} gives, for each partition
$\lambda \vdash n$,
\begin{equation}\label{eq:lower}
  \rk(\Pi|_{V_\lambda}) \;\ge\; \dim(V_\lambda^H).
\end{equation}

Section~\ref{sec:total-rank} establishes (via computational
verification and the image basis conjecture) that the total rank
satisfies
\begin{equation}\label{eq:total}
  \sum_{\lambda \vdash n} \dim(V_\lambda) \cdot \rk(\Pi|_{V_\lambda})
  \;=\; \frac{n!}{2^{\lfloor n/2 \rfloor}}.
\end{equation}

Combining \eqref{eq:lower} and \eqref{eq:total} with
Theorem~\ref{thm:frobenius}:
\[
  0 \;=\;
  \sum_{\lambda} \dim(V_\lambda)
  \bigl(\rk(\Pi|_{V_\lambda}) - \dim(V_\lambda^H)\bigr).
\]
Each summand is a product of $\dim(V_\lambda) > 0$ and
$\rk(\Pi|_{V_\lambda}) - \dim(V_\lambda^H) \ge 0$.
A sum of non-negative terms equaling zero forces each term to
vanish:
\[
  \rk(\Pi|_{V_\lambda}) \;=\; \dim(V_\lambda^H)
  \qquad \text{for all } \lambda \vdash n.
  \qedhere
\]
\end{proof}


%% ====================================================================
\section{Total rank}\label{sec:total-rank}
%% ====================================================================

\begin{proposition}[Total rank formula]\label{prop:total-rank}
In the regular representation of\/ $\Sym_n$,
\[
  \rk(\Pi) \;=\; \frac{n!}{2^{\lfloor n/2 \rfloor}}.
\]
\end{proposition}

\begin{proof}[Computational verification]
Direct computation of $\Pi$ in the regular representation
confirms this formula for $n \le 6$:

\medskip
\begin{center}
\begin{tabular}{ccc}
\toprule
$n$ & $n!$ & $n!/2^{\lfloor n/2 \rfloor}$ \\
\midrule
2 & 2 & 1 \\
3 & 6 & 3 \\
4 & 24 & 6 \\
5 & 120 & 30 \\
6 & 720 & 90 \\
\bottomrule
\end{tabular}
\end{center}
\medskip

In each case the computed rank of $\Pi$ matches the predicted
value exactly.
\end{proof}

\begin{conjecture}[Image basis]\label{conj:image-basis}
A permutation $\sigma \in \Sym_n$ is \emph{descent-paired} if,
for every $k \in \{1, 3, 5, \ldots, 2\lfloor n/2 \rfloor - 1\}$,
we have $\sigma(k) < \sigma(k+1)$ (i.e., position $k$ is not a
descent). The set of descent-paired permutations has cardinality
$n!/2^{\lfloor n/2 \rfloor}$ (OEIS A090932), and their images
$\{\Pi(e_\sigma)\}$ form a basis for $\im(\Pi)$.
\end{conjecture}

\begin{remark}
Conjecture~\ref{conj:image-basis} has been verified for $n \le 5$
(linear independence of the image vectors). It provides a
combinatorial route to a proof of Proposition~\ref{prop:total-rank}:
the count of descent-paired permutations is $n!/2^{\lfloor n/2\rfloor}$
by the Callan characterization, and if their images are linearly
independent, the rank is at least this; if they span the image
(which follows from the orbit structure of $H$ on $\Sym_n$),
the rank is exactly this.
\end{remark}


%% ====================================================================
\section{Computational verification}\label{sec:computation}
%% ====================================================================

All computations were performed in SageMath using exact arithmetic
over $\mathbb{Q}$.

\begin{proposition}[Per-irrep verification]\label{prop:per-irrep}
For all $n \le 8$ and all partitions $\lambda \vdash n$ (a total
of 73 irreducible representations):
\[
  \rk(\Pi|_{V_\lambda}) \;=\; \dim(V_\lambda^H).
\]
\end{proposition}

\begin{proposition}[Step-by-step non-annihilation]\label{prop:step-verify}
For all $n \le 8$ and all $\lambda \vdash n$, the dimension of the
transported image $W_j = P_{i_j} \cdots P_{i_1}(V_\lambda^H)$
remains equal to $\dim(V_\lambda^H)$ at every step $j = 1, \ldots,
\binom{n}{2}$.
\end{proposition}

\begin{proposition}[Total rank verification]\label{prop:total-verify}
For $n \le 6$, the rank of $\Pi$ in the regular representation
equals $n!/2^{\lfloor n/2 \rfloor}$.
\end{proposition}

\begin{proposition}[Image basis verification]\label{prop:basis-verify}
For $n \le 5$, the images $\{\Pi(e_\sigma)\}$ for descent-paired
$\sigma$ are linearly independent and span $\im(\Pi)$.
\end{proposition}

\begin{proposition}[Even-block injectivity]\label{prop:even-verify}
For even block starts $P_{m+1}$ with $m + 1 \le 14$,
the parabolic reduction argument verifies injectivity
unconditionally. This covers all block transitions for $n \le 16$.
\end{proposition}

\medskip

\begin{center}
\begin{tabular}{lll}
\toprule
\textbf{Component} & \textbf{Status} & \textbf{Range} \\
\midrule
Per-irrep $\rk = \dim V^H$ & Verified & $n \le 8$ (73 irreps) \\
Step-by-step non-annihilation & Verified & $n \le 8$ \\
Total rank formula & Verified & $n \le 6$ (regular rep) \\
Image basis conjecture & Verified & $n \le 5$ \\
Even-block parabolic reduction & Verified & $k \le 14$ ($n \le 16$) \\
\midrule
Within-block injectivity & \textbf{Proved} & All $n$ \\
Odd-block boundary injectivity & \textbf{Proved} & All $n$ (contraction) \\
Block $2$ boundary injectivity & \textbf{Proved} & All $n$ (adj.\ disjoint.) \\
Even-block boundary ($k \ge 4$) & \textbf{[Gap]} & General $m$ \\
Total rank (unconditional proof) & \textbf{[Gap]} & General $n$ \\
\bottomrule
\end{tabular}
\end{center}

\bigskip

\begin{remark}[Architecture of the proof]
Three of the four components of Proposition~\ref{prop:injectivity}
are proved unconditionally:
\emph{within-block} (Corollary~\ref{cor:within-block}),
\emph{block~$2$ start} (Lemma~\ref{lem:even-boundary}),
\emph{odd block starts} (Lemma~\ref{lem:odd-boundary}).
Only the even-block boundaries for $k \ge 4$ remain.

The proof is structured so that either remaining gap
(even-block injectivity or total rank formula) would close
the theorem:
\begin{enumerate}
\item If even-block injectivity is proved for all $m$, then
  Proposition~\ref{prop:injectivity} gives
  $\rk \ge \dim V^H$ per irrep, and the Frobenius identity
  forces equality.
\item Alternatively, if the total rank
  $\rk(\Pi) = n!/2^{\lfloor n/2\rfloor}$ is proved independently
  (e.g., via Conjecture~\ref{conj:image-basis}), then combined
  with the per-irrep lower bound, the Frobenius identity
  forces equality.
\end{enumerate}
\end{remark}

\end{document}
